\apendice{Manual del desarrollador e investigador}

Este manual proporciona la información necesaria para comprender la arquitectura interna del software, facilitar su mantenimiento y permitir la replicación del entorno de investigación por parte de otros desarrolladores o personal científico.

\section{Estructura de directorios}
El proyecto se organiza en una estructura de carpetas que separa claramente el entorno de ejecución, la lógica de la aplicación, los scripts de la base de datos y la documentación técnica en LaTeX:

\begin{itemize}
	\item \ruta{/Backend/}: Carpeta principal que contiene la lógica del servidor Flask.
	\begin{itemize}
		\item \ruta{app.py}: Archivo que aloja las rutas y algoritmo de filtrado.
        		\item \ruta{extension.py}: Archivo de configuración donde se realiza la instancia de la base de datos (SQLAlchemy).
       		\item \ruta{modelos.py}: Definición de las clases y tablas de relación (mapeo ORM).
        		\item \ruta{templates/}: Directorio con las plantillas HTML (Jinja2) para la renderización de las vistas.
       		\item \ruta{static/}: Contiene los archivos CSS para el diseño visual y el archivo \texttt{Procfile} para el despliegue en la nube.
        		\item \ruta{.env}: Archivo de configuración de variables de entorno (claves secretas y URLs de conexión). 
	\end{itemize}
    	\item \ruta{/Base\_de\_datos/}: Contiene los scripts SQL necesarios para la inicialización.
    	\begin{itemize}
      		  \item \ruta{TFG.sql}: Script para la creación y población de la base de datos en entorno local.
        		  \item \ruta{TFG\_nube.sql}: Versión optimizada del script para el funcionamiento en servidores de producción.
    	\end{itemize}
   	 \item \ruta{/img/}: Almacena los diagramas y capturas de pantalla utilizados en la memoria y anexos.
    	\item \ruta{/tex/}: Contiene los archivos fuente de LaTeX (\texttt{.tex}) que componen los capítulos y apéndices del trabajo.
   	\item \ruta{/venv/}: Entorno virtual de Python con las librerías aisladas del sistema.
    	\item \ruta{Archivos en la raíz}:
    	\begin{itemize}
        		\item \ruta{requirements.txt}: Lista de requisitos del proyecto para su instalación directa.
        		\item \ruta{gitignore}: Configuración de archivos y carpetas excluidos del almacenamiento en Github.
        		\item \ruta{README.md}: Documentación de presentación del repositorio.
        		\item \ruta{memoria.tex} y \ruta{anexos.tex}: Archivos maestros para la compilación de la documentación final.
   	 \end{itemize}
\end{itemize}

\section{Diseño de clases y Modelado ORM}
La persistencia de la aplicación se basa en un mapeo objeto-relacional (ORM) mediante SQLAlchemy. A continuación, se presenta el diagrama de clases que representa la jerarquía y las relaciones del catálogo SAAC.

\imagen{../img/diagrama_clases.png}{Diagrama de clases UML del modelo de datos.}

Este esquema destaca la centralidad de la clase \texttt{SAACSistema}, la cual actúa como principal para las relaciones de muchos a muchos (\textit{Many-to-Many}) con las tablas maestras, permitiendo que el algoritmo realice consultas cruzadas eficientes durante el filtrado.

\section{Especificación del diccionario de datos (Input Usuario)}
El algoritmo de filtrado por exclusión recibe los datos del formulario en un objeto JSON estructurado. Este diccionario es el que permite realizar el "efecto embudo", descartando sistemas según se comparan las claves del JSON con los requisitos de la base de datos.

\subsection{Estructura del diccionario JSON}
A continuación se detalla la especificación del objeto \texttt{user\_data} procesado en la ruta \texttt{/recomendar}:

\begin{verbatim}
{
    "input_usuario": {
        "visual": int,       // Nivel 0-3 (Criterio de exclusión)
        "auditivo": int,     // Nivel 0-3
        "tecnologico": int,  // Nivel 0-3
        "habla": int,        // Nivel 0-3 (Criterio para bancos de voz)
        "entradas": [int],   // Lista de IDs de la tabla tipo_entrada
        "plataformas": [int],// Lista de IDs de la tabla plataforma
        "idiomas": [int],    // Lista de IDs de la tabla idioma
        "metodos": [int],    // Lista de IDs (Alfabeto/Pictogramas)
        "entorno": int       // ID de entorno (Domicilio/Exterior)
    }
}
\end{verbatim}

\subsection{Lógica de Filtrado por Exclusión (Efecto Embudo)}
El algoritmo opera mediante una cascada de filtros secuenciales que reducen el conjunto total de SAAC ($S_{total}$) hasta obtener la recomendación final ($S_{rec}$):

\begin{enumerate}
    \item \textbf{Filtro de capacidades críticas:} Se eliminan sistemas cuyos \texttt{sistema\_requisito\_funcional} sean superiores a los niveles declarados por el usuario.
    \item \textbf{Filtro de acceso:} Se excluyen sistemas que no soporten ninguno de los \texttt{metodos\_entrada} disponibles para el paciente.
    \item \textbf{Filtro tecnológico:} Se filtran sistemas que no operen en las \texttt{plataformas} accesibles para el usuario.
    \item \textbf{Filtro de entorno:} Si el usuario requiere uso en exterior, se descartan sistemas marcados únicamente como domésticos.
\end{enumerate}

\section{Compilación, instalación y ejecución del proyecto}
\subsection{Entorno de desarrollo}
Se recomienda el uso de \texttt{pip} y \texttt{virtualenv}. El sistema utiliza \texttt{python-dotenv} para cargar la configuración, por lo que es necesario definir \texttt{DATABASE\_URL} y \texttt{FLASK\_ENV} en el entorno local antes de la ejecución.

\subsection{Base de datos y ORM}
La persistencia se gestiona mediante \textbf{SQLAlchemy 2.0}. Para inicializar el esquema en un nuevo entorno de base de datos, se debe ejecutar desde la consola de Python:
\begin{verbatim}
from app import db
with app.app\_context():
	db.create\_all()
\end{verbatim}

\section{Pruebas del sistema}
Para validar la fiabilidad del recomendador y la integridad de los datos, se han realizado los siguientes bloques de pruebas:

\begin{itemize}
	\item \textbf{Pruebas funcionales (Caja negra):} Validación de los formularios de cuestionario y feedback, asegurando que los métodos POST envían correctamente los resultados al servidor.
	\item \textbf{Pruebas de lógica de filtrado:} Introducción de perfiles clínicos extremos (ej. ausencia total de movilidad) para comprobar que el algoritmo filtra correctamente los SAAC incompatibles.
	\item \textbf{Validación de persistencia:} Comprobación en el panel de administración (\texttt{/admin\_historial}) que los datos guardados coinciden exactamente con los enviados por el usuario.
	\item \textbf{Pruebas de seguridad:} Verificación de las cabeceras de seguridad inyectadas por \texttt{Flask-Talisman} (HSTS, XSS Protection) mediante la web VirusTotal \cite{virustotal:upload}.
\end{itemize}

\section{Instrucciones para la modificación o mejora del proyecto}
\subsection{Ampliación del catálogo de SAAC}
Para añadir nuevos sistemas o periféricos no es necesario modificar el código fuente, basta con insertar los nuevos registros en las tablas correspondientes de PostgreSQL, siguiendo el esquema de capacidades requerido por el algoritmo de filtrado.

\subsection{Mejora del algoritmo de recomendación}
 Se recomienda el uso de técnicas de \textit{Machine Learning} en futuras versiones, utilizando el historial de \texttt{feedback} recolectado en este proyecto como conjunto de entrenamiento (Dataset) para refinar las sugerencias automáticas.